% % Reemplace por el título de su capítulo.
\chapter{La pausa}

% % Reemplace por un epígrafe genérico y su atribución.
\epigraph{Detenerse a tiempo ordena la marcha.}{Atribuido a un dicho común}

% % Reemplace por la escena inicial: qué se detiene y por qué.
A mitad de la semana el trabajo pide una pausa y el cuerpo la agradece. Se
cierra el cuaderno, se apaga la lámpara y se sale a caminar sin destino
fijo. El barrio enseña su ritmo lento: una tienda abierta, un perro
dormido y dos vecinos que hablan sin prisa. Entonces queda claro que,
\enquote{parar también es parte del plan}.

% % Reemplace por el método de su pausa: los pasos que sigue.
\section{ Tres pasos para ordenar la pausa}

% % Reemplace por su propio flujo; mantenga el estilo y el título.
\begin{mermaid}[style=monochrome, width=0.85\linewidth, caption={De la idea a la forma final}, pos=H]
  flowchart LR
  A[Idea] --> B[Borrador] --> C[Final]
\end{mermaid}

% % Reemplace por lo que anota en su cuaderno durante la pausa.
La pausa sigue tres pasos simples: anotar lo pendiente, elegir lo primero
y soltar el resto. El diagrama resume ese orden sin adornos.

\begin{center}
  * * *
\end{center}

% % Reemplace por su fragmento breve: un poema o una nota de registro.
% % Reemplace por su fragmento; máximo doce líneas cortas.
\begin{code}[lang=text]
nota del martes, sin prisa:
la mesa queda libre,
la taza vuelve a su lugar,
la lista espera en calma,
la tarde sigue su curso,
mañana será otro día
\end{code}

% % Reemplace por el cierre: qué trae de vuelta la pausa.
Al volver, la mesa está igual pero la cabeza está distinta. Lo urgente
sigue ahí, mas ya no manda. La pausa devuelve el orden y el trabajo
continúa con mejor pulso.
